\subsection{Milestone 1 -- Costruzione del dataset}

Il dataset prodotto contiene le seguenti caratteristiche:

\begin{table}[h!]
\centering
\begin{tabular}{@{}ll@{}}
\toprule
\textbf{Proprietà} & \textbf{Valore} \\
\midrule
Progetto & Apache OpenJPA \\
Versioni Jira totali & 42 \\
Release considerate (primo 33\%) & 14 \\
Release saltate (nessun tag Git) & 3 \\
Snapshot totali & 15.338 \\
Snapshot etichettati come buggy & 1.378 (8,98\%) \\
Commit totali nel repository & 7.811 \\
Commit con chiave ticket & 5.430 (69\%) \\
Ticket unici indicizzati & 1.815 \\
\bottomrule
\end{tabular}
\caption{Riepilogo del dataset prodotto per Apache OpenJPA.}
\end{table}

Le 14 release considerate coprono il periodo da novembre 2006 (\texttt{0.9.6}) ad aprile
2010 (\texttt{2.0.0-beta3}). L'ordinamento cronologico delle release non coincide sempre
con l'ordinamento lessicografico dei nomi, poiché il progetto sviluppava in parallelo
branch di manutenzione e di sviluppo.

Su 15.338 snapshot totali, 1.378 risultano buggy (circa il 9\%). La distribuzione non è
uniforme tra le release. Una possibile spiegazione è che nelle fasi iniziali del progetto
gli sviluppatori compilassero il campo \textit{Affected Versions} su Jira con minore
accuratezza, portando a una sottostima della bugginess nelle release più vecchie: i bug
esistevano, ma i ticket corrispondenti non documentavano esplicitamente le versioni
affette.
